\section{AHSME 1984}
        \begin{enumerate}
        	\item (\textit{AHSME 1984}) $\frac{1000^2}{252^2-248^2}$ equals


 $\mathrm{(A) \  }62500 \qquad \mathrm{(B) \  }1000 \qquad \mathrm{(C) \  } 500\qquad \mathrm{(D) \  }250 \qquad \mathrm{(E) \  } \frac{1}{2}$


 

	\item (\textit{AHSME 1984}) If $x, y$, and $y-\frac{1}{x}$ are not $0$, then


 
 \[\frac{x-\frac{1}{y}}{y-\frac{1}{x}}\] 
equals


 $\mathrm{(A) \ }1 \qquad \mathrm{(B) \ }\frac{x}{y} \qquad \mathrm{(C) \ } \frac{y}{x}\qquad \mathrm{(D) \ }\frac{x}{y}-\frac{y}{x} \qquad \mathrm{(E) \ } xy-\frac{1}{xy}$


 

	\item (\textit{AHSME 1984}) Let $n$ be the smallest nonprime \href{https://artofproblemsolving.com/wiki/index.php?title=Integer}{integer} greater than $1$ with no \href{https://artofproblemsolving.com/wiki/index.php?title=Prime\_factorization}{prime factor} less than $10$. Then 


 $\mathrm{(A) \ }100<n\leq110 \qquad \mathrm{(B) \ }110<n\leq120 \qquad \mathrm{(C) \ } 120<n\leq130 \qquad \mathrm{(D) \ }130<n\leq140 \qquad \mathrm{(E) \ } 140<n\leq150$


 

	\item (\textit{AHSME 1984}) A rectangle intersects a circle as shown: $AB = 4$, $BC = 5$ and $DE = 3$. Then $EF$ equals

                    \begin{figure}[H]
                        \centering
                        \includegraphics[width=0.5\textwidth]{Problem/AHSME_I_1984_4_Problem_Figure_1.jpg}
                    \end{figure}
                    


 $\mathrm{(A) \ }6 \qquad \mathrm{(B) \ }7 \qquad \mathrm{(C) \ }\frac{20}{3} \qquad \mathrm{(D) \ }8 \qquad \mathrm{(E) \ }9$


 

	\item (\textit{AHSME 1984}) The largest \href{https://artofproblemsolving.com/wiki/index.php?title=Integer}{integer} $n$ for which $n^{200}<5^{300}$ is


 $\mathrm{(A) \ }8 \qquad \mathrm{(B) \ }9 \qquad \mathrm{(C) \ } 10 \qquad \mathrm{(D) \ }11 \qquad \mathrm{(E) \ } 12$


 

	\item (\textit{AHSME 1984}) In a certain school, there are three times as many boys as girls and nine times as many girls as teachers. Using the letters $b, g, t$ to represent the number of boys, girls and teachers, respectively, then the total number of boys, girls and teachers can be represented by the \href{https://artofproblemsolving.com/wiki/index.php?title=Expression}{expression}


 $\mathrm{(A) \ }31b \qquad \mathrm{(B) \ }\frac{37}{27}b \qquad \mathrm{(C) \ } 13g \qquad \mathrm{(D) \ }\frac{37}{27}g \qquad \mathrm{(E) \ } \frac{37}{27}t$


 

	\item (\textit{AHSME 1984}) When Dave walks to school, he averages $90$ steps per minute, each of his steps $75$ cm long. It takes him $16$ minutes to get to school. His brother, Jack, going to the same school by the same route, averages $100$ steps per minute, but his steps are only $60$ cm long. How long does it take Jack to get to school?


 $\mathrm{(A) \ }14 \frac{2}{9} \text{ minutes} \qquad \mathrm{(B) \ }15 \text{ minutes}\qquad \mathrm{(C) \ } 18 \text{ minutes}\qquad \mathrm{(D) \ }20 \text{ minutes} \qquad \mathrm{(E) \ } 22 \frac{2}{9} \text{ minutes}$


 

	\item (\textit{AHSME 1984}) Figure $ABCD$ is a \href{https://artofproblemsolving.com/wiki/index.php?title=Trapezoid}{trapezoid} with $AB||DC$, $AB=5$, $BC=3\sqrt{2}$, $\angle BCD=45^\circ$ and $\angle CDA=60^\circ$. The length of $DC$ is


 
                    \begin{figure}[H]
                        \centering
                        \includegraphics[width=0.5\textwidth]{Problem/AHSME_I_1984_8_Problem_Figure_1.jpg}
                    \end{figure}
                    


 $\mathrm{(A) \ }7+\frac{2}{3}\sqrt{3} \qquad \mathrm{(B) \ }8 \qquad \mathrm{(C) \ } 9 \frac{1}{2} \qquad \mathrm{(D) \ }8+\sqrt{3} \qquad \mathrm{(E) \ } 8+3\sqrt{3}$


 

	\item (\textit{AHSME 1984}) The number of \href{https://artofproblemsolving.com/wiki/index.php?title=Digit}{digits} in $4^{16}5^{25}$ (when written in the usual \href{https://artofproblemsolving.com/wiki/index.php?title=Base\_number}{base} $10$ form) is


 $\mathrm{(A) \ }31 \qquad \mathrm{(B) \ }30 \qquad \mathrm{(C) \ } 29 \qquad \mathrm{(D) \ }28 \qquad \mathrm{(E) \ } 27$


 

	\item (\textit{AHSME 1984}) Four \href{https://artofproblemsolving.com/wiki/index.php?title=Complex\_numbers}{complex numbers} lie at the \href{https://artofproblemsolving.com/wiki/index.php?title=Vertices}{vertices} of a \href{https://artofproblemsolving.com/wiki/index.php?title=Square}{square} in the \href{https://artofproblemsolving.com/wiki/index.php?title=Complex\_plane}{complex plane}. Three of the numbers are $1+2i, {-2}+i$ and ${-1}-2i$. The fourth number is 


 $\mathrm{(A) \ }2+i \qquad \mathrm{(B) \ }2-i \qquad \mathrm{(C) \ } 1-2i \qquad \mathrm{(D) \ }{-1}+2i \qquad \mathrm{(E) \ } {-2}-i$


 

	\item (\textit{AHSME 1984}) A calculator has a key that replaces the displayed entry with its \href{https://artofproblemsolving.com/wiki/index.php?title=Perfect\_square}{square}, and another key which replaces the displayed entry with its \href{https://artofproblemsolving.com/wiki/index.php?title=Reciprocal}{reciprocal}. Let $y$ be the final result when one starts with an entry $x\not=0$ and alternately squares and reciprocates $n$ times each. Assuming the calculator is completely accurate (e.g. no roundoff or overflow), then $y$ equals


 $\mathrm{(A) \ }x^{((-2)^n)} \qquad \mathrm{(B) \ }x^{2n} \qquad \mathrm{(C) \ } x^{-2n} \qquad \mathrm{(D) \ }x^{-(2^n)} \qquad \mathrm{(E) \ } x^{((-1)^n2n)}$


 

	\item (\textit{AHSME 1984}) If the \href{https://artofproblemsolving.com/wiki/index.php?title=Sequence}{sequence} $\{a_n\}$ is defined by


 
 \[a_1=2,\]

 \[a_{n+1}=a_n+2n \qquad (n \geq 1),\]


 then $a_{100}$ equals


 $\mathrm{(A) \ }9900 \qquad \mathrm{(B) \ }9902 \qquad \mathrm{(C) \ } 9904 \qquad \mathrm{(D) \ }10100 \qquad \mathrm{(E) \ } 10102$


 

	\item (\textit{AHSME 1984}) $\frac{2\sqrt{6}}{\sqrt{2}+\sqrt{3}+\sqrt{5}}$ equals


 $\mathrm{(A) \ }\sqrt{2}+\sqrt{3}-\sqrt{5} \qquad \mathrm{(B) \ }4-\sqrt{2}-\sqrt{3} \qquad \mathrm{(C) \ } \sqrt{2}+\sqrt{3}+\sqrt{6}-5 \qquad$


 $\mathrm{(D) \ }\frac{1}{2}(\sqrt{2}+\sqrt{5}-\sqrt{3}) \qquad \mathrm{(E) \ } \frac{1}{3}(\sqrt{3}+\sqrt{5}-\sqrt{2})$


 

	\item (\textit{AHSME 1984}) The product of all real \href{https://artofproblemsolving.com/wiki/index.php?title=Root\_(polynomial)}{roots} of the equation $x^{\log_{10}{x}}=10$ is


 $\mathrm{(A) \ }1 \qquad \mathrm{(B) \ }-1 \qquad \mathrm{(C) \ } 10 \qquad \mathrm{(D) \ }10^{-1} \qquad \mathrm{(E) \ } \text{None of these}$


 

	\item (\textit{AHSME 1984}) If $\sin{2x}\sin{3x}=\cos{2x}\cos{3x}$, then one value for $x$ is


 $\mathrm{(A) \ }18^\circ \qquad \mathrm{(B) \ }30^\circ \qquad \mathrm{(C) \ } 36^\circ \qquad \mathrm{(D) \ }45^\circ \qquad \mathrm{(E) \ } 60^\circ$


 

	\item (\textit{AHSME 1984}) The \href{https://artofproblemsolving.com/wiki/index.php?title=Function}{function} $f(x)$ satisfies $f(2+x)=f(2-x)$ for all real numbers $x$. If the equation $f(x)=0$ has exactly four distinct real \href{https://artofproblemsolving.com/wiki/index.php?title=Root\_(polynomial)}{roots}, then the sum of these roots is


 $\mathrm{(A) \ }0 \qquad \mathrm{(B) \ }2 \qquad \mathrm{(C) \ } 4 \qquad \mathrm{(D) \ }6 \qquad \mathrm{(E) \ } 8$


 

	\item (\textit{AHSME 1984}) A \href{https://artofproblemsolving.com/wiki/index.php?title=Right\_triangle}{right triangle} $ABC$ with \href{https://artofproblemsolving.com/wiki/index.php?title=Hypotenuse}{hypotenuse} $AB$ has side $AC=15$. \href{https://artofproblemsolving.com/wiki/index.php?title=Altitude}{Altitude} $CH$ divides $AB$ into segments $AH$ and $HB$, with $HB=16$. The \href{https://artofproblemsolving.com/wiki/index.php?title=Area}{area} of $\triangle ABC$ is


 
                    \begin{figure}[H]
                        \centering
                        \includegraphics[width=0.5\textwidth]{Problem/AHSME_I_1984_17_Problem_Figure_1.jpg}
                    \end{figure}
                    


 $\mathrm{(A) \ }120 \qquad \mathrm{(B) \ }144 \qquad \mathrm{(C) \ } 150 \qquad \mathrm{(D) \ }216 \qquad \mathrm{(E) \ } 144\sqrt{5}$


 

	\item (\textit{AHSME 1984}) A point $(x, y)$ is to be chosen in the \href{https://artofproblemsolving.com/wiki/index.php?title=Coordinate\_plane}{coordinate plane} so that it is equally distant from the $x$-axis, the $y$-axis, and the \href{https://artofproblemsolving.com/wiki/index.php?title=Line}{line} $x+y=2$. Then $x$ is


 $\mathrm{(A) \ }\sqrt{2}-1 \qquad \mathrm{(B) \ }\frac{1}{2} \qquad \mathrm{(C) \ } 2-\sqrt{2} \qquad \mathrm{(D) \ }1 \qquad \mathrm{(E) \ } \text{not uniquely determined}$


 

	\item (\textit{AHSME 1984}) A box contains $11$ balls, numbered $1, 2, 3,\ldots,11$. If $6$ balls are drawn simultaneously at random, what is the \href{https://artofproblemsolving.com/wiki/index.php?title=Probability}{probability} that the sum of the numbers on the balls drawn is odd?


 $\mathrm{(A) \ }\frac{100}{231} \qquad \mathrm{(B) \ }\frac{115}{231} \qquad \mathrm{(C) \ } \frac{1}{2} \qquad \mathrm{(D) \ }\frac{118}{231} \qquad \mathrm{(E) \ } \frac{6}{11}$


 

	\item (\textit{AHSME 1984}) The number of the distinct solutions of the \href{https://artofproblemsolving.com/wiki/index.php?title=Equation}{equation} $|{x-|2x+1|}|=3$ is


 $\mathrm{(A) \ }0 \qquad \mathrm{(B) \ }1 \qquad \mathrm{(C) \ } 2 \qquad \mathrm{(D) \ }3 \qquad \mathrm{(E) \ } 4$


 

	\item (\textit{AHSME 1984}) The number of triples $(a, b, c)$ of \href{https://artofproblemsolving.com/wiki/index.php?title=Natural\_numbers}{positive integers} which satisfy the simultaneous \href{https://artofproblemsolving.com/wiki/index.php?title=Equation}{equations}


 
 \[ab+bc=44,\]

 \[ac+bc=23,\]


 is


 $\mathrm{(A) \ }0 \qquad \mathrm{(B) \ }1 \qquad \mathrm{(C) \ } 2 \qquad \mathrm{(D) \ }3 \qquad \mathrm{(E) \ } 4$


 

	\item (\textit{AHSME 1984}) Let $a$ and $c$ be fixed \href{https://artofproblemsolving.com/wiki/index.php?title=Natural\_numbers}{positive numbers}. For each \href{https://artofproblemsolving.com/wiki/index.php?title=Real\_number}{real number} $t$ let $(x_t, y_t)$ be the \href{https://artofproblemsolving.com/wiki/index.php?title=Vertex}{vertex} of the \href{https://artofproblemsolving.com/wiki/index.php?title=Parabola}{parabola} $y=ax^2+tx+c$. If the set of the vertices $(x_t, y_t)$ for all real values of $t$ is graphed on the \href{https://artofproblemsolving.com/wiki/index.php?title=Cartesian\_plane}{plane}, the \href{https://artofproblemsolving.com/wiki/index.php?title=Graph}{graph} is


 $\mathrm{(A) \ } \text{a straight line} \qquad \mathrm{(B) \ } \text{a parabola} \qquad \mathrm{(C) \ } \text{part, but not all, of a parabola} \qquad \mathrm{(D) \ } \text{one branch of a hyperbola} \qquad \mathrm{(E) \ } \text{none of these}$


 

	\item (\textit{AHSME 1984}) $\frac{\sin{10^\circ}+\sin{20^\circ}}{\cos{10^\circ}+\cos{20^\circ}}$ equals


 $\mathrm{(A) \ }\tan{10^\circ}+\tan{20^\circ} \qquad \mathrm{(B) \ }\tan{30^\circ} \qquad \mathrm{(C) \ } \frac{1}{2}(\tan{10^\circ}+\tan{20^\circ}) \qquad \mathrm{(D) \ }\tan{15^\circ} \qquad \mathrm{(E) \ } \frac{1}{4}\tan{60^\circ}$


 

	\item (\textit{AHSME 1984}) If $a$ and $b$ are positive \href{https://artofproblemsolving.com/wiki/index.php?title=Real\_numbers}{real numbers} and each of the \href{https://artofproblemsolving.com/wiki/index.php?title=Equation}{equations} $x^2+ax+2b=0$ and $x^2+2bx+a=0$ has real \href{https://artofproblemsolving.com/wiki/index.php?title=Root\_(polynomial)}{roots}, then the smallest possible value of $a+b$ is


 $\mathrm{(A) \ }2 \qquad \mathrm{(B) \ }3 \qquad \mathrm{(C) \ } 4 \qquad \mathrm{(D) \ }5 \qquad \mathrm{(E) \ } 6$


 

	\item (\textit{AHSME 1984}) The total \href{https://artofproblemsolving.com/wiki/index.php?title=Area}{area} of all the \href{https://artofproblemsolving.com/wiki/index.php?title=Faces}{faces} of a \href{https://artofproblemsolving.com/wiki/index.php?title=Rectangular\_prism}{rectangular solid} is $22 \ \text{cm}^2$, and the total length of all its \href{https://artofproblemsolving.com/wiki/index.php?title=Edges}{edges} is $24 \ \text{cm}$. Then the length in cm of any one of its \href{https://artofproblemsolving.com/wiki/index.php?title=Diagonal\#Polyhedra}{interior diagonals} is


 $\mathrm{(A) \ }\sqrt{11} \qquad \mathrm{(B) \ }\sqrt{12} \qquad \mathrm{(C) \ } \sqrt{13}\qquad \mathrm{(D) \ }\sqrt{14} \qquad \mathrm{(E) \ } \text{not uniquely determined}$


 

	\item (\textit{AHSME 1984}) In the \href{https://artofproblemsolving.com/wiki/index.php?title=Obtuse\_triangle}{obtuse triangle} $ABC$, $AM=MB$, $MD\perp BC$, $EC\perp BC$. If the \href{https://artofproblemsolving.com/wiki/index.php?title=Area}{area} of $\triangle ABC$ is $24$, then the area of $\triangle BED$ is


 
                    \begin{figure}[H]
                        \centering
                        \includegraphics[width=0.5\textwidth]{Problem/AHSME_I_1984_26_Problem_Figure_1.jpg}
                    \end{figure}
                    


 $\mathrm{(A) \ }9 \qquad \mathrm{(B) \ }12 \qquad \mathrm{(C) \ } 15 \qquad \mathrm{(D) \ }18 \qquad \mathrm{(E) \ } \text{not uniquely determined}$


 

	\item (\textit{AHSME 1984}) In $\triangle ABC$, $D$ is on $AC$ and $F$ is on $BC$. Also, $AB\perp AC$, $AF\perp BC$, and $BD=DC=FC=1$. Find $AC$.


 
                    \begin{figure}[H]
                        \centering
                        \includegraphics[width=0.5\textwidth]{Problem/AHSME_I_1984_27_Problem_Figure_1.jpg}
                    \end{figure}
                    


 $\mathrm{(A) \ }\sqrt{2} \qquad \mathrm{(B) \ }\sqrt{3} \qquad \mathrm{(C) \ } \sqrt[3]{2} \qquad \mathrm{(D) \ }\sqrt[3]{3} \qquad \mathrm{(E) \ } \sqrt[4]{3}$


 

	\item (\textit{AHSME 1984}) The number of distinct pairs of \href{https://artofproblemsolving.com/wiki/index.php?title=Integers}{integers} $(x, y)$ such that $0<x<y$ and $\sqrt{1984}=\sqrt{x}+\sqrt{y}$ is


 $\mathrm{(A) \ }0 \qquad \mathrm{(B) \ }1 \qquad \mathrm{(C) \ } 3 \qquad \mathrm{(D) \ }4 \qquad \mathrm{(E) \ } 7$


 

	\item (\textit{AHSME 1984}) Find the largest value for $\frac{y}{x}$ for pairs of \href{https://artofproblemsolving.com/wiki/index.php?title=Real\_numbers}{real numbers} $(x, y)$ which satisfy $(x-3)^2+(y-3)^2=6$.


 $\mathrm{(A) \ }3+2\sqrt{2} \qquad \mathrm{(B) \ }2+\sqrt{3} \qquad \mathrm{(C) \ } 3\sqrt{3} \qquad \mathrm{(D) \ }6 \qquad \mathrm{(E) \ } 6+2\sqrt{3}$


 

	\item (\textit{AHSME 1984}) For any \href{https://artofproblemsolving.com/wiki/index.php?title=Complex\_number}{complex number} $w=a+bi$, $|w|$ is defined to be the \href{https://artofproblemsolving.com/wiki/index.php?title=Real\_number}{real number} $\sqrt{a^2+b^2}$. If $w=\cos40^\circ+i\sin40^\circ$, then
$|w+2w^2+3w^3+...+9w^9|^{-1}$ equals


 $\mathrm{(A) \ }\frac{1}{9}\sin40^\circ \qquad \mathrm{(B) \ }\frac{2}{9}\sin20^\circ \qquad \mathrm{(C) \ } \frac{1}{9}\cos40^\circ \qquad \mathrm{(D) \ }\frac{1}{18}\cos20^\circ \qquad \mathrm{(E) \ } \text{none of these}$


 

\end{enumerate}